% ---------------------------------------------------------------------------
%  00-how-to-read.tex -- orientation for the two audiences this document has
% ---------------------------------------------------------------------------
\section*{How to read this document}
\addcontentsline{toc}{section}{How to read this document}

This document has two readers in mind at once, and it tries not to compromise
for either.

\begin{itemize}[leftmargin=1.2em, itemsep=2pt]
\item \textbf{The reader who does not do timed automata.}  Every section opens
      with the idea in ordinary language, and wherever the technical text gets
      dense there is an \emph{In plain words} box restating the same content
      with no notation.  You can read only the opening paragraphs, the boxes,
      and the figure captions and come away with an accurate picture.

\item \textbf{The reader who will referee this.}  The formal statements,
      proofs, and the exact boundary between ``ours'' and ``standard'' are in
      \cref{sec:theory} and \cref{sec:novelty}.  \cref{sec:novelty} is
      deliberately blunt: it lists what is new, what is a port of somebody
      else's idea, and what is engineering that carries no theoretical claim.
      If you only read one section, read that one.
\end{itemize}

\paragraph{Everything is traced to code.}
Each mechanism names the file and the function that implements it, so any
claim here can be checked against the implementation rather than taken on
trust:

\begin{center}
\begin{tabular}{@{}ll@{}}
\toprule
\textbf{Mechanism} & \textbf{Where it lives} \\
\midrule
level-synchronous parallel search & \file{src/netreach.cpp}, \code{reach\_engine} \\
striped passed set, two-phase lock & \file{src/netreach.cpp}, \code{NetPassed}, \code{SpinRW} \\
antichain update (subsume + merge)& \file{src/netreach.cpp}, \code{NetPassed::add} \\
lock-free coverage probe & \file{src/netreach.cpp}, \code{NetPassed::probe\_kill} \\
zone interner (64 shards) & \file{src/netreach.cpp}, \code{ZoneIntern} \\
zone algebra, $\aLU$, convex union & \file{HDBM/include/ptctlmc/dbm\_ops.hpp} \\
minimal-graph codec and inclusion & \file{HDBM/include/ptctlmc/mindbm.hpp} \\
per-thread zone allocator & \file{HDBM/include/ptctlmc/dbm\_pool.hpp} \\
\bottomrule
\end{tabular}
\end{center}

\paragraph{A note on honesty.}
Several things in the implementation look clever and are not: they were tried,
measured, and are load-bearing only because a measurement said so.  Several
other things were tried and \emph{removed}, and those are recorded too
(\cref{sec:limits}).  A TACAS submission is stronger for naming the ideas that
did not work than for pretending the design was obvious.
